\fullname

The quiz contains 9 questions. Questions 8 and 9 are worth 2 points each, all others 1 point. The total is 11 points and the score is out of 10 (so you still get max points if you lose one point).

One handwritten A6 (or 4x6 index card) cheat sheet allowed, no printed material.
Remember that your electronic devices should be silent and put away in your bag at all times.

\bigskip
\begin{center}\rule{0.4\linewidth}{0.4pt}\end{center}
\bigskip

\begin{question}
\lstinline|Deque<E>| declares \lstinline|addFirst|, \lstinline|addLast|, \lstinline|removeFirst|, \lstinline|removeLast| (among others). Suppose you only ever add elements with \textbf{one} of these two methods, and only ever remove them with \textbf{one} of the other two (always the same pair). For each pair below, mark whether it makes the deque behave like a \lstinline|Stack|, a \lstinline|Queue|, or neither.

\begin{center}
\begin{tabular}{@{}p{2.6in}ccc@{}}
\toprule
Add with / remove with & Stack & Queue & Neither \\
\midrule
\lstinline|addFirst| / \lstinline|removeFirst| & \cb{1} & \cb{0} & \cb{0} \\
\lstinline|addLast| / \lstinline|removeLast| & \cb{1} & \cb{0} & \cb{0} \\
\lstinline|addLast| / \lstinline|removeFirst| & \cb{0} & \cb{1} & \cb{0} \\
\lstinline|addFirst| / \lstinline|removeLast| & \cb{0} & \cb{1} & \cb{0} \\
\bottomrule
\end{tabular}
\end{center}
\end{question}

\begin{question}
Consider the following code, where \lstinline|N| is a positive integer.

\begin{lstlisting}
/* q is defined as below */
for (int i = 0; i < N; i++)
    q.offer(i);
while (!q.isEmpty())
    System.out.println(q.poll());
\end{lstlisting}

Which definition of \lstinline|q| gives the lowest worst-case number of operations for this code (select one)? Then write that complexity as a function of $N$.

\begin{itemize}[leftmargin=2em,itemsep=1pt,topsep=2pt]
\item[\cb{0}] \lstinline|Queue<Integer> q = new ArrayDeque<>();|
\item[\cb{0}] \lstinline|Queue<Integer> q = new LinkedList<>();|
\item[\cb{1}] both have the same complexity
\end{itemize}

Complexity: $O($\underline{\hspace{1in}}$)$ \ifanswers ($N$)\fi
\end{question}

\begin{question}
Consider the following declarations.

\begin{lstlisting}
Deque<Integer> d = new LinkedList<>();
List<Integer> l = new LinkedList<>();
\end{lstlisting}

Mark the statements that compile (ignore what happens when they run).
\begin{itemize}[leftmargin=2em,itemsep=1pt,topsep=2pt]
\item[\cb{1}] \lstinline|d.push(1);|
\item[\cb{0}] \lstinline|d.get(0);|
\item[\cb{1}] \lstinline|l.get(0);|
\item[\cb{0}] \lstinline|l.push(1);|
\end{itemize}
\end{question}

\begin{question}
The following method of an \lstinline|ArrayQueue| (circular array, fields \lstinline|data|, \lstinline|front|, \lstinline|size|) removes and returns the front element, but one line contains a bug.
Which change fixes the method (select one)?

Hint: try it by hand on a queue where \lstinline|front| is not 0 and the queue wraps around the end of the array.

\begin{lstlisting}[numbers=left,numberstyle=\tiny,xleftmargin=2em]
public E dequeue() {
    if (size == 0) throw new NoSuchElementException("empty queue");
    E e = data[front];
    data[front] = null;
    front = front + 1;
    size--;
    return e;
}
\end{lstlisting}

\begin{itemize}[leftmargin=2em,itemsep=1pt,topsep=2pt]
\item[\cb{0}] line 3 becomes \lstinline|E e = data[front % data.length];|
\item[\cb{0}] line 4 becomes \lstinline|data[front % data.length] = null;|
\item[\cb{1}] line 5 becomes \lstinline[keepspaces=true]|front = (front + 1) % data.length;|
\item[\cb{0}] line 5 becomes \lstinline[keepspaces=true]|front = (front + 1) % size;|
\item[\cb{0}] line 6 becomes \lstinline[keepspaces=true]|size = size - 1 % data.length;|
\end{itemize}
\end{question}

\begin{question}
The following four lines are the body of \lstinline|addFirst(e)| for a circular-array deque (field \lstinline|data|, \lstinline|front|, \lstinline|size|; method \lstinline|grow()| doubles the capacity), but they are in the wrong order.
Write the numbers 1 to 4 in the boxes to give the order in which they must be executed.

\medskip
\begin{tabular}{@{}l@{\quad}l@{}}
\ansbox{4} & \lstinline|size++;| \\[3pt]
\ansbox{1} & \lstinline[keepspaces=true]|if (size == data.length) grow();| \\[3pt]
\ansbox{3} & \lstinline|data[front] = e;| \\[3pt]
\ansbox{2} & \lstinline[keepspaces=true]|front = (front - 1 + data.length) % data.length;| \\
\end{tabular}
\end{question}

\begin{question}
What does the following code fragment print?

\begin{lstlisting}
Queue<Integer> q1 = new ArrayDeque<>();
Queue<Integer> q2 = new ArrayDeque<>();
Queue<Integer> tmp;

q2.offer(40);
while (!q1.isEmpty()) q2.offer(q1.poll());
tmp = q1; q1 = q2; q2 = tmp;    // swap q1 and q2

q2.offer(15);
while (!q1.isEmpty()) q2.offer(q1.poll());
tmp = q1; q1 = q2; q2 = tmp;    // swap q1 and q2

q2.offer(70);
while (!q1.isEmpty()) q2.offer(q1.poll());
tmp = q1; q1 = q2; q2 = tmp;    // swap q1 and q2

System.out.print(q1.poll() + " " + q1.poll() + " " + q1.poll());
\end{lstlisting}

Answer: \underline{\hspace{2in}} \ifanswers (70 15 40: the two queues emulate a stack)\fi
\end{question}

\begin{question}
\lstinline|dq| is an \lstinline|ArrayDeque<Integer>| with $N$ elements, and \lstinline|k| is a non-negative \lstinline|int|.

\begin{lstlisting}
void rotate(ArrayDeque<Integer> dq, int k) {
    int steps = k % dq.size();
    for (int i = 0; i < steps; i++)
        dq.addLast(dq.removeFirst());
}
\end{lstlisting}

\lstinline|rotate| moves the first \lstinline|steps| elements to the back, preserving their relative order. What is the worst-case complexity of \lstinline|rotate|, as a function of $N$?

Answer: \underline{\hspace{1.5in}} \ifanswers ($O(N)$: \lstinline|steps| $< N$ for any \lstinline|k|)\fi
\end{question}

\begin{question}
(2 points) A standard deque implementation of \lstinline|Integer|s, based on a circular array, starts empty with an initial capacity of 4 (its capacity doubles when it grows). The following operations are performed in order:

\begin{minipage}[t]{0.3\linewidth}
\vspace{-\baselineskip}
\begin{lstlisting}
dq.addFirst(1);
dq.addLast(2);
dq.addFirst(3);
dq.removeLast();
dq.addLast(4);
dq.addFirst(5);
dq.removeFirst();
\end{lstlisting}
\end{minipage}\hfill
\begin{minipage}[t]{0.66\linewidth}
Draw the final state of the underlying array (every cell, whether it holds a reference to an \lstinline|Integer| object or none at all), the \lstinline|Integer| objects it references, and the final values of \lstinline|front| and \lstinline|size|.

Hint: the cells that do not reference any object contain \lstinline|null|.
\end{minipage}

\begin{answerspace}{2.5in}
\ifanswers
\begin{verbatim}
array (cells 0..3): [ref->4, null, ref->3, ref->1]
front = 2, size = 3
(the deque never holds more than 4 elements, so the array never grows;
2 is removed by removeLast(), 5 by removeFirst())
\end{verbatim}
\fi
\end{answerspace}
\end{question}

\begin{question}
(2 points) Design a class \lstinline|RecentLog| that keeps only the last \lstinline|k| lines of a log, where \lstinline|k| is given to the constructor \lstinline|RecentLog(int k)|. It supports two methods, \lstinline|add| and \lstinline|get|, that should run in $O(1)$ time in the \textbf{worst case}, regardless of \lstinline|k|.
\begin{itemize}[leftmargin=2em,itemsep=1pt,topsep=2pt]
\item \lstinline|add(String line)|: adds \lstinline|line|; if \lstinline|k| lines are already stored, the oldest one is discarded;
\item \lstinline|get(int i)|: returns the \lstinline|i|-th most recent line stored (\lstinline|get(0)| is the line added last); throws an exception if \lstinline|i| is negative or not smaller than the number of lines stored.
\end{itemize}

Your answer should briefly outline what fields you would have in your data structure (their type and how they are initialized), how \lstinline|add| and \lstinline|get| work, and why they achieve the intended complexity. You do not need to write code.

\begin{answerspace}{4.5in}
\ifanswers
\begin{verbatim}
Fields: String[] data of length k (allocated once, never grows),
        int front (index of the oldest line), int size.

add(line): if size < k: data[(front + size) % k] = line; size++
           else (full): data[front] = line (overwrites the oldest);
                        front = (front + 1) % k
           O(1) worst case: no shifting, and no grow() since the
           capacity is fixed.

get(i):    throw if i < 0 or i >= size
           return data[(front + size - 1 - i) % k]
           O(1): one index computation.

Also correct (simpler): String[] data of length k, long n = number
of lines ever added.
add(line): data[n % k] = line; n++
get(i):    throw if i < 0 or i >= min(n, k)
           return data[(n - 1 - i) % k]

Not O(1) worst case: ArrayDeque/LinkedList have no O(1) get(i);
a growing array is only O(1) amortized for add.
\end{verbatim}
\fi
\end{answerspace}
\end{question}
